../cictikz/src/cictikz/data/tex/cictikz_preamble.tex

% What a binary code actually does on the way from 1 to 2, drawn twice.
%
% Left, MSB first: 01 -> 11 -> 10, so the output goes 1, 3, 2 and overshoots.
% Right, LSB first: 01 -> 00 -> 10, so it goes 1, 0, 2 and undershoots.
% Either way the code passes through a value that is not between the two it
% is moving between, which is the glitch the lecture is after. The four
% states are drawn in both panels, unvisited ones included, so the two
% routes can be read against the full state space of dac_bin_states.

\begin{tikzpicture}[thick]

  ../cictikz/src/cictikz/data/tex/dacsm_lib.tex

  % ---- MSB first ----
  \dacState{a00}{(0,3)}{00}
  \dacState{a01}{(3.6,3)}{01}
  \dacState{a11}{(0,0)}{11}
  \dacState{a10}{(3.6,0)}{10}
  \dacArrow{a01}{a11}
  \dacArrow{a11}{a10}

  % ---- LSB first ----
  \dacState{b00}{(7.6,3)}{00}
  \dacState{b01}{(11.2,3)}{01}
  \dacState{b11}{(7.6,0)}{11}
  \dacState{b10}{(11.2,0)}{10}
  \dacArrow{b01}{b00}
  \dacArrow{b00}{b10}

\end{tikzpicture}
\end{document}
